\documentclass[12pt]{article}
\usepackage[T1]{fontenc}
\usepackage[utf8]{inputenc}
\usepackage[polish]{babel}
\usepackage{lmodern}
\usepackage[a4paper,margin=2.4cm]{geometry}
\usepackage{xcolor}
\usepackage{graphicx}
\usepackage{amsmath,amssymb}
\usepackage{booktabs}
\usepackage{array}
\usepackage{enumitem}
\usepackage{hyperref}
\usepackage{tcolorbox}
\usepackage{fancyhdr}
\definecolor{WEIIPrimary}{RGB}{0,130,150}
\definecolor{WEIIPrimaryDark}{RGB}{0,86,100}
\definecolor{WEIIPrimaryLight}{RGB}{229,246,248}
\definecolor{PRzNavy}{RGB}{31,63,142}
\definecolor{RLTextDark}{RGB}{34,34,34}
\definecolor{RLGray}{RGB}{90,90,90}
\definecolor{RLLightGray}{RGB}{245,246,248}
\definecolor{RLWarning}{RGB}{184,86,0}
\definecolor{RLExample}{RGB}{52,112,64}
\definecolor{RLCodeBg}{RGB}{248,248,248}

\usepackage{listings}
\usepackage{xcolor}

\lstdefinestyle{rlpython}{%
    language=Python,
    basicstyle=\ttfamily\scriptsize,
    keywordstyle=\color{WEIIPrimaryDark}\bfseries,
    commentstyle=\color{RLGray}\itshape,
    stringstyle=\color{RLExample},
    numbers=left,
    numberstyle=\tiny\color{RLGray},
    numbersep=6pt,
    showstringspaces=false,
    breaklines=true,
    tabsize=4,
    frame=single,
    rulecolor=\color{WEIIPrimary!35},
    backgroundcolor=\color{RLCodeBg},
    captionpos=b
}

\lstset{style=rlpython}


\hypersetup{colorlinks=true,linkcolor=WEIIPrimaryDark,urlcolor=WEIIPrimaryDark,citecolor=WEIIPrimaryDark}
\setlength{\headheight}{15pt}
\pagestyle{fancy}
\fancyhf{}
\lhead{Uczenie się ze wzmocnieniem}
\rhead{WEiI PRz}
\cfoot{\thepage}

\newtcolorbox{taskbox}[1]{colback=WEIIPrimaryLight,colframe=WEIIPrimary,title=\textbf{#1}}
\newtcolorbox{notebox}[1]{colback=RLLightGray,colframe=RLGray,title=\textbf{#1}}
\newtcolorbox{warninghandout}[1]{colback=RLWarning!8,colframe=RLWarning,title=\textbf{#1}}

\newcommand{\handouttitle}[3]{%
\begin{center}
{\Large\bfseries #1}\\[0.2cm]
{\large #2}\\[0.15cm]
{\small #3}
\end{center}
\vspace{0.4cm}
\hrule
\vspace{0.6cm}
}

\newcommand{\EnableLabNumbering}{%
  \renewcommand{\thesection}{L\arabic{section}}%
  \renewcommand{\thesubsection}{\thesection.\arabic{subsection}}%
  \renewcommand{\thesubsubsection}{\thesubsection.\arabic{subsubsection}}%
}


\title{Raport projektowy\\Uczenie się ze wzmocnieniem}
\author{Joanna Joniec}
\date{2026/2027}

\begin{document}
\maketitle
\tableofcontents
\newpage

\section{Informacje podstawowe}

\begin{center}
\begin{tabular}{p{0.34\textwidth}p{0.56\textwidth}}
\toprule
Element & Informacja \\
\midrule
Imię i nazwisko & Joanna Joniec \\
Numer albumu & 177089 \\
Kierunek/specjalność & Informatyka/ Sztuczna Inteligencja \\
Tytuł projektu & "Planowanie trasy kuriera w dynamicznym środowisku miejskim z wykorzystaniem uczenia ze wzmocnieniem" \\
Link do repozytorium & \url{https://github.com/JoAsha2004/uczenie-sie-ze-wzmocnieniem---PROJEKT} \\
Data oddania & 15.11.2026 \\
\bottomrule
\end{tabular}
\end{center}

\section{Temat i cel projektu}

\subsection{Opis problemu}
Projekt ma na celu symulację trasy kuriera rozwożącego pizzę na dwuwymiarowej planszy podzielonej na pola reprezentujące uproszczoną sieć miejskich ulic.
Punktem startowym dla agenta poruszającego się pomiędzy sąsiednimi polami planszy (każdy ruch wiąże się z upływem czasu i zużyciem określonej ilości energii) jest pizzeria, z której odbiera określoną liczbę zamówień, a punktami docelowymi są znajdujący się w różnych, oddalonych od siebie lokalizacjach mapy klienci, do których musi dotrzeć przed powrotem do pizzerii po kolejne przesyłki (nowy epizod).
Agent omija trwale nieprzejezdne pola (budynki) oraz bramy, które są początkowo otwarte, ale mogą zostać zamknięte do końca obecnego epizodu, jeśli kurier przeszedł przez dane wyznaczone pole określoną wcześniej liczbę razy w czasie obecnego epizodu.
Im lokalizacja dostawy znajduje się dalej od pizzerii, tym oferuje wyższy początkowy napiwek doliczany do stałej zapłaty, ale jego wartość zmniejsza się krokowo wraz z wydłużającym się czasem dostawy, co wymusza decyzję o kolejności obsługi klientów. 
Jeśli napiwek spadnie poniżej 0, klient anuluje zamówienie i kurier już go nie dostarcza (co skutkuje mniejszym zarobkiem na koniec epizodu).
W przypadku gdy energia spadła do zera agent przenosi się do punktu startowego i rozpoczyna się kolejny epizod.

\subsection{Cel agenta}
Zadaniem agenta jest zaplanowanie efektywną trasę od pizzerii do poszczególnych klientów oraz najkrótszą trasę powrotną od ostatniego klienta (co kończy epizod) do pizzerii, nie znając optymalnego rozwiązania, ale ucząc się na podstawie otrzymywanych nagród i kar. 
Agent musi dostarczyć jak najwięcej zamówień i powrócić do pizzerii przed wyczerpaniem energii.
Kurier dąży do uzyskania jak najwyższej łącznej nagrody, czyli sumy uzyskanych dochodów (stałe zapłaty + napiwki) pomniejszoną o połowę potencjalnego zarobku z utraconych zamówień. 

\subsection{Uzasadnienie zastosowania uczenia się ze wzmocnieniem}
Uczenie się ze wzmocnieniem jest uzasadnione w przypadku tego projektu, ponieważ w dynamicznym środowisku gridowym nie ma wcześniej ustalonej optymalnej drogi, ale istnieje wiele potencjalnych tras, którymi agent może poruszać się aby osiągnąć wcześniej ustalony cel, a uczenie się przez eksplorację środowiska, sukcesy i porażki umożliwia znajdywanie jak najefektywniejszych tras w kolejnych epizodach. Najbliższy klient lub najwyższy bieżący napiwek nie musi prowadzić do największego zysku na końcu epizodu.
Aby dotrzeć do celów agent podejmuje szereg kolejnych decyzji o kierunku ruchu zależnych od jego aktualnego położenia względem przeszkód i klientów.
Stan środowiska zmienia się w zależności od działań agenta. Przejazd przez określone pola zwiększa licznik wykorzystania przejścia, co przy zbyt często powtarzających się trasach doprowadza do permanentnego do końca epizodu zamknięcia przejazdu. Każdy ruch wiąże się ze spadkiem energii a także upływem czasu. Agent decyduje długoterminowo, które dostawy są priorytetowe - przyniosą największy dochód przy najmniejszym spadku energii, ale ich wybór wpływa na wydłużony czas oczekiwania pozostałych klientów, co skutkuje zmniejszeniem napiwku a nawet anulowaniem zamówienia. 


\section{Formalizacja problemu RL/MDP}

\subsection{Agent i środowisko}
W tym projekcie agentem jest kurier, który sam wybiera kolejne ruchy nie znając optymalnej strategii i uczy się poprzez interakcję ze środowiskiem. 
Środowiskiem dla agenta jest plansza gridowa reprezentująca miasto, na której znajdują się pizzeria (punkt początkowy i końcowy), klienci (których musi odwiedzić, aby dostać zarobek i napiwek), przeszkody (budynki i zamykające się bramy) oraz ograniczony zasób energii. 
Agent i środowisko oddziałują ze sobą w kolejnych krokach epizodu. Agent, obserwując aktualną sytuację, wybiera kolejny ruch, z kolei środowisko aktualizuje pozycję agenta, po czym zmniejsza jego energię, aktualizuje czas oczekiwania klientów, sprawdza warunki zamknięcia bram, zmienia stan i oblicza nagrodę.
W pojedynczym kroku epizodu agent najpierw otrzymuje aktualny stan, potem wybiera akcję, a środowisko próbuje wykonać ruch, po czym pozycja agenta jest aktualizowana, energia się zmniejsza, zwiększa się czas dostawy, a napiwki klientów są aktualizowane. Następnie sprawdzane są warunki zamknięcia bram, nagroda jest wyznaczana i jest generowany nowy stan.

\subsection{Stan lub obserwacja}
\begin{center}
\begin{tabular}{p{0.30\textwidth}p{0.60\textwidth}}
\toprule
Składowa stanu/obserwacji & Znaczenie \\
\midrule
Pozycja kuriera (x,y) & Określa aktualne położenie agenta na planszy i umożliwia wybór kierunku dalszego ruchu \\
Aktualna energia & Informuje o liczbie pozostałych jednostek energii dostępnych do wykonania dalszych ruchów \\
Status zamówień (dostarczone / aktywne / anulowane) & Pozwala określić, które zamówienia zostały już dostarczone, a które nadal wymagają realizacji \\
Stan bram (otwarta / zamknięta) & Określa dostępność poszczególnych przejazdów i wpływa na możliwość wyboru trasy \\
Aktualna wartość napiwków & Informuje o potencjalnym zysku możliwym do uzyskania od poszczególnych klientów w danym momencie \\
Licznik przejazdów przez bramę & Pozwala przewidzieć możliwość przyszłego zamknięcia przejazdu po przekroczeniu ustalonego limitu użyć \\
\bottomrule
\end{tabular}
\end{center}

\subsection{Akcje}
\begin{center}
\begin{tabular}{p{0.30\textwidth}p{0.60\textwidth}}
\toprule
Akcja & Opis \\
\midrule
Góra & Agent porusza się o 1 kratkę w górę \\
Dół & Agent porusza się o 1 kratkę w dół \\
Lewo & Agent porusza się o 1 kratkę w lewo \\
Prawo & Agent porusza się o 1 kratkę w prawo \\
\bottomrule
\end{tabular}
\end{center}
W przypadku gdy akcja powodowałaby wejście w budynek, zamkniętą bramę lub poza granice mapy, ruch jest niedozwolony, a agent pozostaje na aktualnym polu i traci energię jak w przypadku wykonania właściwego ruchu.

\subsection{Funkcja nagrody}
\begin{equation}
r_t = \begin{cases}
+20 & \text{dostarczenie zamówienia} \\
+tip_i & \text{uzyskanie napiwku od klienta} \\
-0.1 & \text{wykonanie ruchu} \\
-2 & \text{próba wejścia w przeszkodę} \\
-10 & \text{utrata zamówienia} \\
-20 & \text{wyczerpanie energii przed zakończeniem epizodu}
\end{cases}
\end{equation}

Funkcja nagrody została zaprojektowana w taki sposób, aby premiować realizację głównego celu zadania, czyli skuteczne dostarczanie zamówień przy jednoczesnym ograniczaniu czasu przejazdu i zużycia energii. Funkcja zawiera zarówno nagrody dodatnie związane z realizacją dostaw, jak i kary za działania utrudniające osiągnięcie celu. Dobrze zaprojektowana funkcja nagrody powinna premiować oczekiwany efekt działania, a nie konkretną strategię, którą agent ma stosować.
\\Agent otrzymuje dodatnią nagrodę za dostarczenie zamówienia do klienta. Jest to główny sygnał sukcesu w środowisku. Dodatkowo agent otrzymuje premię równą aktualnej wartości napiwku klienta. Dzięki temu opłacalne staje się nie tylko dostarczenie zamówienia, ale również wykonanie dostawy możliwie szybko. Takie rozwiązanie zachęca agenta do planowania kolejności obsługi klientów oraz uwzględniania przyszłych konsekwencji własnych decyzji.
\\Funkcja nagrody zawiera również kary. Każdy wykonany ruch powoduje naliczenie niewielkiej kary −0.1, pełniącej rolę kosztu kroku. Kara za krok zachęca agenta do wybierania krótszych i bardziej efektywnych tras oraz ogranicza ryzyko przypadkowego poruszania się po planszy bez realizowania celu. Wartość tej kary została dobrana tak, aby była istotnie mniejsza od nagrody za dostawę zamówienia i nie zniechęcała agenta do eksploracji środowiska.Dodatkowa kara −2 jest przyznawana za próbę wejścia w budynek, zamkniętą bramę lub poza granice planszy. Takie działanie nie zmienia pozycji agenta, ale powoduje stratę energii i pogarsza końcowy wynik epizodu. Kara ta ma zniechęcać do wykonywania niedozwolonych akcji.
\\W środowisku występują zarówno ograniczenia twarde, jak i miękkie. Do ograniczeń twardych należą niemożność wejścia na pole zajęte przez budynek lub zamkniętą bramę oraz zakończenie epizodu po wyczerpaniu energii. Ograniczeniami miękkimi są natomiast koszt każdego ruchu, spadek wartości napiwków wraz z upływem czasu oraz utrata części potencjalnego zysku wynikająca z anulowania zamówienia.
\\Przyjęta funkcja nagrody nie powinna zachęcać agenta do niepożądanego zachowania. Największe dodatnie nagrody są związane wyłącznie z realizacją dostaw, natomiast zwlekanie, wykonywanie zbędnych ruchów lub ignorowanie klientów prowadzi do stopniowego obniżania wyniku. Oznacza to, że maksymalizacja sumy nagród jest zgodna z rzeczywistym celem środowiska, którym jest obsłużenie możliwie największej liczby klientów przy zachowaniu wysokiego dochodu i racjonalnego zużycia energii. Należy jednak podkreślić, że przedstawione wartości nagród i kar mają charakter wstępny i mogą zostać skorygowane po przeprowadzeniu pierwszych eksperymentów.

\subsection{Warunki zakończenia epizodu}

Epizod rozpoczyna się w momencie opuszczenia pizzerii przez kuriera z określoną liczbą zamówień i kończy się po spełnieniu jednego z poniższych warunków:
\begin{itemize} \item agent dostarczył wszystkie aktywne zamówienia i powrócił do pizzerii, \item poziom energii agenta spadł do zera, \item wszystkie pozostałe zamówienia zostały anulowane przez klientów wskutek przekroczenia maksymalnego czasu oczekiwania, \item została osiągnięta maksymalna liczba kroków dozwolonych w pojedynczym epizodzie. \end{itemize}
Pierwszy przypadek traktowany jest jako zakończenie sukcesem. Pozostałe sytuacje oznaczają zakończenie niepowodzeniem lub częściową realizację zadania. Wprowadzenie limitu energii oraz maksymalnej liczby kroków zapobiega powstawaniu bardzo długich epizodów, w których agent porusza się po planszy bez realizacji celu. Jasne zdefiniowanie warunków zakończenia epizodu jest niezbędnym elementem formalizacji problemu RL.

\subsection{Założenia i ograniczenia środowiska} Przyjęto, że środowisko ma postać dwuwymiarowej planszy gridowej o skończonych wymiarach. Agent porusza się wyłącznie pomiędzy sąsiednimi polami mapy, wykonując ruch w jednym z czterech kierunków. Środowisko jest deterministyczne, co oznacza, że wykonanie tej samej akcji w tym samym stanie zawsze prowadzi do tego samego rezultatu. W celu uproszczenia modelu przyjęto, że wszyscy klienci są znani agentowi od początku epizodu, a ich lokalizacja nie zmienia się w trakcie jego trwania. Nie uwzględniono również ruchu innych pojazdów, zmiennych warunków drogowych, pogody ani losowych zdarzeń mogących wpływać na czas przejazdu. Budynki stanowią przeszkody statyczne i nie mogą zostać przekroczone przez agenta. Dodatkowo w środowisku występują dynamiczne przeszkody w postaci bram, które mogą zostać zamknięte po przekroczeniu określonego limitu wykorzystania przejazdu. Zamknięcie bramy obowiązuje wyłącznie do końca bieżącego epizodu. Przyjęto, że każda wykonana akcja powoduje identyczny koszt energetyczny oraz jednakowy przyrost czasu niezależnie od kierunku ruchu. Wartość napiwków zmniejsza się wraz z czasem oczekiwania klientów zgodnie z ustalonymi regułami środowiska. Model środowiska stanowi uproszczoną reprezentację rzeczywistego problemu logistycznego. Celem projektu nie jest wierne odwzorowanie procesu dostawy pizzy, lecz stworzenie środowiska umożliwiającego analizę zachowania agenta RL w zadaniu wymagającym planowania trasy, zarządzania ograniczonymi zasobami oraz przewidywania przyszłych konsekwencji podejmowanych decyzji.

\section{Implementacja}

\subsection{Struktura repozytorium}

Należy krótko opisać strukturę repozytorium i najważniejsze pliki projektu.

\begin{verbatim}
repozytorium/
|-- src/
|-- experiments/
|-- report/
|-- results/
`-- README.md
\end{verbatim}

\subsection{Środowisko}

Należy opisać sposób implementacji środowiska, w szczególności funkcje inicjalizacji epizodu oraz wykonania kroku środowiska.

\subsection{Agent bazowy}

Należy opisać agenta losowego, heurystycznego albo inny baseline wykorzystany do porównania z agentem uczącym się.

\subsection{Agent uczący się}

Należy opisać zastosowany algorytm RL lub DRL, jego parametry oraz sposób implementacji. W przypadku Q-tab należy opisać reprezentację tablicy wartości. W przypadku DRL należy opisać wejścia modelu, wyjścia modelu, architekturę oraz sposób uczenia.

\section{Plan eksperymentów}

Należy opisać eksperymenty zaplanowane w celu oceny działania agenta.

\begin{center}
\begin{tabular}{p{0.24\textwidth}p{0.30\textwidth}p{0.30\textwidth}}
\toprule
Eksperyment & Co jest zmieniane? & Co jest mierzone? \\
\midrule
E1 & Do uzupełnienia & Do uzupełnienia \\
E2 & Do uzupełnienia & Do uzupełnienia \\
E3 & Do uzupełnienia & Do uzupełnienia \\
\bottomrule
\end{tabular}
\end{center}

W tej sekcji należy wskazać:
\begin{itemize}
    \item liczbę epizodów uczenia,
    \item liczbę powtórzeń eksperymentu albo uzasadnienie braku powtórzeń,
    \item parametry uczenia,
    \item metryki oceny,
    \item warianty środowiska, map lub funkcji nagrody,
    \item sposób porównania z agentem bazowym.
\end{itemize}

\section{Wyniki}

W tej sekcji należy przedstawić wyniki eksperymentów. Wyniki powinny być pokazane w postaci tabel, wykresów lub czytelnych zestawień.

\begin{center}
\begin{tabular}{p{0.28\textwidth}p{0.20\textwidth}p{0.20\textwidth}p{0.20\textwidth}}
\toprule
Wariant & Średni zwrot & Skuteczność & Długość epizodu \\
\midrule
Agent bazowy & Do uzupełnienia & Do uzupełnienia & Do uzupełnienia \\
Agent RL/DRL & Do uzupełnienia & Do uzupełnienia & Do uzupełnienia \\
\bottomrule
\end{tabular}
\end{center}

\begin{figure}[h]
    \centering
    \fbox{\parbox{0.82\textwidth}{\centering Miejsce na wykres procesu uczenia}}
    \caption{Przykładowy wykres procesu uczenia agenta.}
\end{figure}

\section{Analiza i wnioski}

Należy zinterpretować wyniki. W tej części należy odpowiedzieć m.in. na pytania:
\begin{itemize}
    \item Czy agent uczący się działa lepiej niż agent bazowy?
    \item Jak zmieniało się zachowanie agenta w trakcie uczenia?
    \item Które elementy środowiska były najtrudniejsze?
    \item Czy funkcja nagrody prowadziła do oczekiwanego zachowania?
    \item Czy wystąpiły przypadki niepożądanego zachowania agenta?
    \item Jakie są najważniejsze ograniczenia rozwiązania?
\end{itemize}

\section{Instrukcja uruchomienia}

Należy opisać sposób uruchomienia projektu od czystego repozytorium. Przykładowo:

\begin{verbatim}
git clone https://github.com/...
cd nazwa-repozytorium
python -m venv .venv
.venv\Scripts\activate
pip install -r requirements.txt
python src/main.py
\end{verbatim}

Należy również wskazać, jak odtworzyć eksperymenty oraz gdzie zapisywane są wyniki.

\section{Wsparcie AI}

W tej sekcji należy obowiązkowo opisać wykorzystanie narzędzi sztucznej inteligencji. Jeżeli student nie korzystał z takich narzędzi, należy wpisać jednoznaczną informację, że nie korzystano ze wsparcia AI.

\begin{center}
\begin{tabular}{p{0.13\textwidth}p{0.19\textwidth}p{0.21\textwidth}p{0.27\textwidth}}
\toprule
Data & Narzędzie/model & Cel użycia & Transkrypt / link / zakres \\
\midrule
Do uzupełnienia & Do uzupełnienia & Do uzupełnienia & Do uzupełnienia \\
Do uzupełnienia & Do uzupełnienia & Do uzupełnienia & Do uzupełnienia \\
\bottomrule
\end{tabular}
\end{center}

Należy zapisać transkrypty rozmów z modelami AI wraz z datą oraz nazwą modelu. Jeżeli rozmowa odbywała się poza narzędziami programistycznymi, należy dodatkowo udostępnić całą rozmowę poprzez link, o ile użyte narzędzie umożliwia takie udostępnienie.

Student powinien także wyjaśnić, które elementy wygenerowane lub zasugerowane przez AI zostały samodzielnie sprawdzone, zmodyfikowane albo odrzucone.

\section{Literatura}

W tej sekcji należy wskazać wykorzystane źródła: książki, artykuły naukowe, dokumentację bibliotek, materiały wykładowe oraz inne źródła użyte podczas pracy nad projektem.

\begin{thebibliography}{9}
\bibitem{sutton_barto}
R. S. Sutton, A. G. Barto, \textit{Reinforcement Learning: An Introduction}, 2nd ed., MIT Press, 2018.
\end{thebibliography}

\end{document}
